% !PRINTED 368
\weberpage{368}
%<<CONT
stens zwei von den $x_3, \dots x_7$ ein Paar darin bilden. Das ist aber nicht möglich, da dieses Paar sonst mit einem der übrigen $x$ ein syzygetisches Tripel bilden würde.

Um nachzuweisen, dass die Doppeltangenten $x_1, x_2, x_3, x_4, x_5, x_6$ immer in einem Complexe vorkommen, nehmen wir zwei beliebige von ihnen, $x_1 x_2$, heraus und wählen ein in dem Complexe $x_1 x_2$ vorkommendes anderes Paar $y_1 y_2$, was auf fünf Arten möglich ist. Daraus bilden wir das syzygetische Complextripel

\weq{2}{\begin{aligned}
\alpha)\ & x_1 x_2,\ y_1 y_2\\
\beta)\ & x_1 y_1,\ x_2 y_2\\
\gamma)\ & x_1 y_2,\ x_2 y_1,
\end{aligned}}

welches fünf verschiedene solcher Tripel repräsentirt. In $\beta$) und $\gamma$) müssen nun $x_3, x_4, x_5, x_6, x_7$ vorkommen, und zwar keine zwei von ihnen gepaart. Wir zeigen nun zunächst, dass sich diese fünf $x$ nicht zu zwei und drei auf die beiden Complexe $\beta$), $\gamma$) vertheilen können, sondern nur zu eins und vier. Nehmen wir nämlich an, die Complexe $\beta$), $\gamma$) seien so zusammengesetzt:

\weq{2}{\begin{aligned}
\beta)\ & x_1 y_1,\ x_2 y_2,\ x_3 y_3,\ x_4 y_4,\ x_5 y_5,\ \cdot\\
\gamma)\ & x_1 y_2,\ x_2 y_1,\ x_6 y_6,\ x_7 y_7,\ \cdot\quad\cdot\text{,}
\end{aligned}}

so können wir noch den Complex

\weq{2}{\delta)\quad x_6 x_7,\ y_6 y_7\ \cdots{},}

betrachten, der mit $\beta$) zusammen ein syzygetisches Paar bildet, weil weder $x_1$ noch $y_1$ in $\delta$) vorkommt, das Paar also nicht azygetisch sein kann. Es müssen also $\beta$) und $\delta$) ein syzygetisches Quadrupel gemein haben, und da in zwei Paaren von $\beta$) mindestens ein von $x_6, x_7$ verschiedenes $x$ vorkommt, so muss dieses $x$ auch in $\delta$) vorkommen, was unmöglich ist, weil kein $x$ mit $x_6, x_7$ syzygetisch ist. Es bleibt also für $\beta$), $\gamma$) nur eine Zusammensetzung übrig, wie die folgende:

\weq{3}{\begin{aligned}
\beta)\ & x_1 y_1,\ x_2 y_2,\ x_3 y_3,\ x_4 y_4,\ x_5 y_5,\ x_6 y_6\\
\gamma)\ & x_1 y_2,\ x_2 y_1,\ x_7 y_7,\ \cdot\quad\cdot\quad\cdot
\end{aligned}}

Dass wir gerade diese Annahme machen, und nicht die sechs $x$ in $\gamma$) aufnehmen, ist keine Beschränkung, da wir, wenn n\"othig, $y_1$ mit $y_2$ vertauschen können.

Wenn wir nun unter Festhaltung von $x_1 x_2$ an Stelle von $y_1 y_2$ die anderen Paare von (2) $\alpha$) treten lassen, so bekommen
%CONT>>
